\section{Methods}

\subsection{Boundary}

This paper treats compression as a finite receipt grammar inside the G900
apparatus. The term gravity shadow is project shorthand for a compression
shadow. It is not a claim about physical gravity, mass, energy, spacetime
curvature, or force generation.

All tests preserve the G900 body. No artifact in this sequence mutates the
body, admits a physical force, or installs a force renderer.

\subsection{Compression metrics}

Compression was measured using finite support metrics only:

\begin{itemize}
  \item support size,
  \item support radius,
  \item mean distance to receipt home,
  \item maximum distance to receipt home,
  \item support entropy,
  \item concentration ratio.
\end{itemize}

These metrics are descriptive. A positive metric result is not sufficient to
admit a global compression shadow.

\subsection{Source discovery}

Candidate compression surfaces were first discovered by scanning inherited
project artifacts and G900 Observatory appliance artifacts. The inherited
project scan covered Projects 19, 29, 30, and 31. The appliance scan covered
the live G900 Observatory source surface.

The discovery and priority stages did not interpret source content. They only
identified candidate files for later inspection.

\subsection{Return-cell transport metric probe}

The first metric-ready surface was the admitted return-cell one-step
information transport. Its support history exposed source slots, target slots,
and a finite target-home proxy. The target slots of the admitted target packet
were used as a conservative receipt-home proxy.

For each source and target slot, distance to the nearest receipt-home slot was
computed in the cyclic 15-slot register.

\subsection{Replay audit}

Replay stability was tested by re-evaluating the return-cell compression
signature under equivalent presentations of the same transport history:

\begin{itemize}
  \item original row order,
  \item reversed row order,
  \item sorted source-target pairs,
  \item unique support reduction.
\end{itemize}

\subsection{Perturbation audit}

Perturbation stability was tested by applying bounded changes to the source
slots, target slots, receipt-home slots, and row support. The audit recorded
which perturbations preserved the compression signature and which failed.

\subsection{Six-nine return-home surface probe}

The six-nine return-home ledger was inspected as a second compression-relevant
surface. It was found to be receipt-strong but not metric-ready as directed
transport. Its rowsets expose surface slotsets rather than an explicit
source-to-target slot history.

\subsection{Surface taxonomy}

The current taxonomy distinguishes two compression-relevant surface types:

\begin{itemize}
  \item return-cell transport, a metric-ready directed transport surface;
  \item six-nine return-home, a receipt-strong surface rowset candidate.
\end{itemize}
